% chapters/03_methoden.tex

\chapter{Methodik}
\label{chap:methoden}

Dieses Kapitel beschreibt das methodische Vorgehen: die Infrastruktur, mit der
das Training reproduzierbar gemacht wurde, den Datensatz und seine Aufbereitung,
das eigentliche Trainingsverfahren mit seinen Hyperparametern, die
Multi-GPU-Skalierung sowie die Methodik der Evaluation. Reproduktionsdetails
(konkrete Kommandos, vollständige Variablenlisten) sind bewusst ausgelagert und
über \doc{docs/}-Verweise zugänglich.

\section{Vorgehensmodell}
\label{sec:vorgehen}

Das Projekt folgte einer klaren Abfolge (\cref{fig:system-overview}): Zunächst
wurde eine containerisierte Trainingsinfrastruktur aufgebaut, anschließend der
Datensatz beschafft und ins erwartete Format konvertiert, danach das Fine-Tuning
auf dem HPC-Cluster durchgeführt und schließlich der resultierende Checkpoint
evaluiert. Die Evaluation -- insbesondere die geschlossene Simulationsschleife --
führte zu Erkenntnissen über den Domain-Gap, die wiederum
Sim-Kalibrierungs- und Konzeptarbeiten anstießen (\cref{chap:ergebnisse}).

\begin{figure}[H]
  \centering
  \resizebox{\linewidth}{!}{%
  \begin{tikzpicture}[
    node distance = 1.1cm and 1.5cm,
    box/.style    = {draw, rounded corners, fill=NavyBlue!8,
                     minimum width=2.9cm, minimum height=1cm,
                     font=\small\sffamily, align=center},
    eval/.style   = {draw, rounded corners, fill=BurntOrange!10,
                     minimum width=2.9cm, minimum height=1cm,
                     font=\small\sffamily, align=center},
    arrow/.style  = {-{Stealth[scale=1.1]}, thick, NavyBlue!70},
  ]
    \node[box]  (data)     {Datensatz\\LeRobot v3.0};
    \node[box]  (convert)  [right=of data]   {Konvertierung\\v3.0\,$\to$\,v2.1};
    \node[box]  (train)    [right=of convert]{Fine-Tuning\\(KISSKI, A100)};
    \node[box]  (ckpt)     [right=of train]  {Checkpoint\\\texttt{checkpoint-*}};
    \node[eval] (openloop) [below=of train]  {Open-Loop-Eval\\(Action-\ac{mse})};
    \node[eval] (closed)   [below=of ckpt]   {Closed-Loop-Eval\\(Isaac~Lab)};

    \draw[arrow] (data)    -- (convert);
    \draw[arrow] (convert) -- (train);
    \draw[arrow] (train)   -- (ckpt);
    \draw[arrow] (ckpt)    -- (closed);
    \draw[arrow] (ckpt.south) |- (openloop.east);
  \end{tikzpicture}}
  \caption{Vorgehensmodell: von der Datenaufbereitung über das Training bis zur
    zweigleisigen Evaluation (Open-Loop und Closed-Loop).}
  \label{fig:system-overview}
\end{figure}

\section{Trainingsinfrastruktur}
\label{sec:infrastruktur}

Eine Kernanforderung war \textbf{Reproduzierbarkeit über heterogene Plattformen}
hinweg: dieselbe Software sollte lokal, auf Cloud-GPUs und auf dem KISSKI-Cluster
identisch laufen. Umgesetzt wurde dies durch einen \textbf{selbst-enthaltenen
Container} (CUDA~12.8, Python~3.10, Paketmanager \texttt{uv}), der die gesamte
Pipeline -- Download, Konvertierung, Training -- über ein Entrypoint-Skript
autonom orchestriert. Lokal läuft das Image unter Docker; auf KISSKI wird es zu
einem Apptainer-SIF konvertiert. Der Trainingscode erkennt beide Umgebungen
selbst, sodass keine Codeänderung nötig ist. Die vollständige Bedienung ist in
\doc{docs/training/anleitung.md} dokumentiert.

\paragraph{Speichermodell.}
Eine bewusste Designentscheidung betrifft die Datenhaltung: Standardmäßig gibt es
\emph{keinen} persistenten Host-Speicher; alle Daten, Checkpoints und Logs liegen
im Container-Dateisystem unter \texttt{/data}. Der Container ist langlebig
gedacht -- \texttt{stop}/\texttt{start} erhält den Zustand, nur \texttt{docker rm}
zerstört ihn. Diese Wahl spiegelt die Semantik von Cloud-GPU-Plattformen wider
(eine Instanz~$\equiv$~ein Container) und verhindert insbesondere den
verbreiteten Fehler, mit \texttt{docker run -{}-rm} unbeabsichtigt
Trainingsergebnisse zu löschen. \textbf{Ausnahme KISSKI:} Dort bindet Apptainer
das VAST-Projektverzeichnis \texttt{/mnt/vast-kisski/\dots/data} nach
\texttt{/data}, sodass Daten über Job-Läufe hinweg persistieren.

\paragraph{Konfiguration über Umgebungsvariablen.}
Sämtliches Verhalten wird über Umgebungsvariablen gesteuert; ihre Defaults sind
im Dockerfile hinterlegt. Dadurch lässt sich derselbe Container ohne Codeänderung
an die jeweilige Hardware anpassen. \Cref{tab:env-curated} zeigt die
konzeptionell wichtigsten Variablen; die vollständige Referenz mit allen Defaults
findet sich in \doc{docs/training/env-vars.md}.

\begin{table}[H]
  \centering
  \begin{tabularx}{\linewidth}{llX}
    \toprule
    \textbf{Variable} & \textbf{Default} & \textbf{Zweck} \\
    \midrule
    \texttt{HF\_TOKEN}          & ---     & Pflicht: HuggingFace-Zugang \\
    \texttt{GLOBAL\_BATCH\_SIZE}& \texttt{8} & globale Batch-Größe (VRAM-abhängig) \\
    \texttt{MAX\_STEPS}         & \texttt{20000} & Trainingsschritte \\
    \texttt{NUM\_GPUS}          & \texttt{1} & GPUs (aktiviert \texttt{torchrun}) \\
    \texttt{LEARNING\_RATE}     & \texttt{1e-4} & Lernrate \\
    \texttt{TUNE\_VISUAL}       & \texttt{0} & Vision-Encoder mittrainieren \\
    \texttt{USE\_AUGMENTATION}  & \texttt{1} & Bild-Augmentierung (Color-Jitter) \\
    \texttt{TRAIN\_TEST\_SPLIT} & \texttt{0} & 80/20-Split aktivieren \\
    \texttt{USE\_COTRAIN}       & \texttt{0} & Co-Training auf echten \emph{und}
      gerenderten Bildern (\cref{sec:cotrain}) \\
    \texttt{SKIP\_*}            & \texttt{0} & einzelne Phasen überspringen
      (\texttt{DOWNLOAD}, \texttt{CONVERT}, \texttt{TRAIN}) \\
    \bottomrule
  \end{tabularx}
  \caption{Auswahl der zentralen Konfigurationsvariablen (Defaults). Vollständige
    Liste: \doc{docs/training/env-vars.md}.}
  \label{tab:env-curated}
\end{table}

\paragraph{Bedienbarkeit als Reproduzierbarkeitsproblem.}
Die Konfiguration über Umgebungsvariablen erkauft ihre Flexibilität mit einer
wachsenden Bedienoberfläche: Im Projektverlauf verteilte sich dieselbe
Parameterliste auf vier Orte -- den Vorgabewert im Skript, dessen
Hilfetext, die Tabelle in der Dokumentation und die Projektübersicht. Solche
Kopien laufen erfahrungsgemäß auseinander, und eine Anleitung, die einen
falschen Vorgabewert nennt, beschädigt genau die Reproduzierbarkeit, um derer
willen die Infrastruktur gebaut wurde.

Die Abhilfe war, die Parameter \emph{einmal} deklarativ zu beschreiben und
alles Weitere daraus abzuleiten: Aus derselben Spezifikation speisen sich ein
geführtes Auswahlmenü für die Bedienung am Terminal, der Hilfetext und eine
Prüfung, die Abweichungen zwischen Spezifikation, Skript und Dokumentation
meldet. Ergänzt wird dies durch einen einzigen Einstiegspunkt, der zunächst
nach der Domäne (Simulation, Training, Cluster) fragt und dann an den
zuständigen Starter übergibt. Zugleich wurden die verbliebenen
maschinengebundenen Pfade entfernt: Der Simulationsserver liest seine
lokalen Einstellungen aus einer nicht versionierten Konfigurationsdatei, die
Cluster-Skripte leiten sämtliche Pfade aus einer einzigen Variablen ab. Damit
läuft das Repository ohne Anpassung auf einer fremden Maschine -- die
praktische Voraussetzung dafür, dass die hier berichteten Läufe von Dritten
wiederholt werden können (\doc{docs/portabilitaet.md}).

\section{Datensatz, Konvertierung und Train/Test-Split}
\label{sec:datensatz}

\paragraph{Datensatz.}
Verwendet wurde der öffentliche Datensatz
\texttt{G1\_Dex3\_BlockStacking}~\cite{g1dex3dataset}: \textbf{301~Episoden}
teleoperierter Block-Stacking-Demonstrationen mit insgesamt \textbf{281\,196}
Bildern bei \qty{30}{\fps}, eine einzige Task. Jeder Zeitschritt umfasst vier
RGB-Kameras (\texttt{cam\_left\_high}, \texttt{cam\_right\_high},
\texttt{cam\_left\_wrist}, \texttt{cam\_right\_wrist}), einen 28-dimensionalen
Propriozeptionsvektor und die Sprach-Task. \Cref{tab:datensatz} fasst die
Kennzahlen zusammen.

\begin{table}[H]
  \centering
  \begin{tabularx}{\linewidth}{lX}
    \toprule
    \textbf{Eigenschaft}       & \textbf{Wert} \\
    \midrule
    Format (Quelle / Ziel)     & LeRobot v3.0 $\to$ v2.1 \\
    Episoden (gesamt)          & 301 \\
    Episoden (Training / Test) & 240 / 61 (80/20-Split) \\
    Bilder gesamt              & 281\,196 \\
    Kameras                    & 4 $\times$ RGB ($640\times480$\,px) \\
    Aufzeichnungsrate          & \qty{30}{\fps} \\
    Zustands-/Aktionsdim.      & 28 (Arme 14, Hände 14) \\
    Aktions-Horizont           & 16 Zeitschritte \\
    \bottomrule
  \end{tabularx}
  \caption{Eigenschaften des Trainings-Datensatzes.}
  \label{tab:datensatz}
\end{table}

\paragraph{Aktionsrepräsentation.}
Der 28-dimensionale Aktionsvektor ist heterogen aufgebaut: Die 14~Armgelenke
(Indizes 0--13) sind als \emph{relative} Gelenk-Deltas kodiert, die 14~Fingergelenke
der DEX3-Hände (Indizes 14--27) als \emph{absolute} Zielwinkel. Diese Mischung
ist später für die korrekte Ansteuerung im Simulator relevant
(\cref{sec:eval-methodik}). Die vollständige Indexbelegung dokumentiert
\doc{docs/training/trainingsverfahren.md}.

\paragraph{Konvertierung.}
Da GR00T das LeRobot-Format~v2.1 erwartet, der Datensatz aber als v3.0 vorliegt,
konvertiert die Pipeline automatisch von v3.0 nach v2.1 und ergänzt bei Bedarf die
\texttt{modality.json}. Dieser Schritt läuft als Teil des Entrypoints und kann
bei bereits konvertierten Daten übersprungen werden.

\paragraph{Train/Test-Split.}
Um eine objektive Bewertung zu ermöglichen, wurde ein \textbf{reproduzierbarer
80/20-Split} konzipiert: die ersten 240~Episoden (Index 0--239) als Trainings-,
die letzten 61 (Index 240--300) als Testmenge.\footnote{Die Grenze liegt bei~240,
nicht bei~241: Die Split-Logik rechnet $n_{\text{train}} = \lfloor 301 \cdot 0{,}8
\rfloor = 240$. Der erste reale Split-Lauf (\cref{sec:lauf3}) protokolliert
entsprechend \texttt{split\_range: "240:301"}.} Bewusst wurde \emph{nicht} zufällig
gesplittet, sondern entlang des natürlichen Episodenindex -- das garantiert
Reproduzierbarkeit ohne Seed-Abhängigkeit. Das Vorgehen und seine
Einschränkungen (ein niedriger Action-Loss auf der Testmenge ist nur ein
\emph{schwacher} Proxy für Aufgabenerfolg) sind in
\doc{docs/training/train-test-split.md} ausgeführt. Dass dieser Split im ersten
durchgeführten Lauf noch \emph{nicht aktiv} war, ist ein Befund der Durchführung
und wird in \cref{sec:split-befund} behandelt.

\section{Trainingsverfahren und Hyperparameter}
\label{sec:trainingsverfahren}

Das Training ist ein \acf{sft} (Post-Training) des vortrainierten Modells per
\acl{bc} -- es ist \emph{kein} Training von Grund auf, kein RL und kein
Multi-Embodiment-Training. Entscheidend ist, dass nur ein Teil des Modells
angepasst wird (\emph{selektives} Fine-Tuning, kein LoRA): Das Sprach-Backbone und
der Vision-Encoder bleiben \textbf{eingefroren}, trainiert werden der Projektor und
der Diffusions-/Flow-Matching-Aktionskopf (\cref{tab:tuning}). Dieses Einfrieren
des Vision-Encoders ist die unmittelbare Ursache des später analysierten
Domain-Gaps -- der Encoder sieht im gesamten Training nie ein Simulationsbild.

\begin{table}[H]
  \centering
  \begin{tabular}{lll}
    \toprule
    \textbf{Komponente} & \textbf{Flag} & \textbf{Status} \\
    \midrule
    Sprach-Backbone (LLM) & \texttt{tune\_llm=False}            & eingefroren \\
    Vision-Encoder        & \texttt{tune\_visual=False}         & eingefroren \\
    Projektor             & \texttt{tune\_projector=True}       & trainiert \\
    Aktionskopf (Diffusion)& \texttt{tune\_diffusion\_model=True}& trainiert \\
    \bottomrule
  \end{tabular}
  \caption{Selektives Fine-Tuning: trainierbare vs.\ eingefrorene Komponenten
    (Standardkonfiguration).}
  \label{tab:tuning}
\end{table}

Die Hyperparameter (\cref{tab:hparams}) wurden für zwei Konfigurationen
festgelegt: einen Einzel-GPU-Lauf (A100, Batch~8, 175\,000~Schritte) und einen
Vier-GPU-Lauf (4$\times$A100, Batch~32, 44\,000~Schritte). Beide verarbeiten rund
1,4~Mio.\ Samples (etwa fünf Epochen) und sind damit äquivalent: $44\,000 \times 32
= 1{,}408\,\text{Mio.} \approx 176\,000 \times 8$. Der erste vollständige Lauf
wurde in der Einzel-GPU-Konfiguration durchgeführt.

\begin{table}[H]
  \centering
  \begin{tabular}{lrr}
    \toprule
    \textbf{Parameter} & \textbf{1$\times$A100} & \textbf{4$\times$A100} \\
    \midrule
    \texttt{GLOBAL\_BATCH\_SIZE}   & 8        & 32 \\
    \texttt{MAX\_STEPS}            & 175\,000 & 44\,000 \\
    \texttt{LEARNING\_RATE}        & 1e-4     & 2e-4 \\
    \texttt{WARMUP\_RATIO}         & 0,05     & 0,05 \\
    \texttt{WEIGHT\_DECAY}         & 1e-5     & 1e-5 \\
    Optimierer / Precision         & \multicolumn{2}{c}{AdamW, bf16 + tf32} \\
    LR-Scheduler                   & \multicolumn{2}{c}{cosine} \\
    Aktions-Horizont               & \multicolumn{2}{c}{16 Schritte} \\
    \bottomrule
  \end{tabular}
  \caption{Trainings-Hyperparameter der beiden Konfigurationen.}
  \label{tab:hparams}
\end{table}

\paragraph{Bild-Augmentierung.}
Standardmäßig aktiv ist eine Augmentierung der Trainingsbilder (im Wesentlichen
Color-Jitter in Helligkeit, Kontrast, Sättigung und Farbton). Sie ist die
abgeschwächte Form der in \cref{sec:domain-gap-grundlagen} genannten
\acf{dr}: Sie variiert das Erscheinungsbild der realen Bilder, ersetzt aber
nicht das Sehen \emph{echter} Simulationsbilder. Ihre Rolle im Projekt ist
dadurch begrenzt, aber nicht unerheblich -- der dritte Lauf
(\cref{sec:lauf3}) ist der erste, der sie zusammen mit dem mittrainierten
Vision-Encoder einsetzt, und unterscheidet sich genau darin vom zweiten.

Die Lernrate des Vier-GPU-Laufs ist mit $\sqrt{4}$ skaliert ($1\text{e-}4 \cdot
\sqrt{4} = 2\text{e-}4$), um dem vierfach größeren Batch Rechnung zu tragen. Eine
Validierung \emph{während} des Trainings findet nicht statt
(\texttt{eval\_strategy="no"}); die Modellqualität wird stattdessen nachgelagert
über die Evaluationsmethodik aus \cref{sec:eval-methodik} bewertet.

\section{Multi-GPU-Skalierung}
\label{sec:multi-gpu}

Da der GR00T-Code die verteilte Trainingslogik (DeepSpeed ZeRO-2,
Gradienten-Synchronisation, rang-bewusstes Logging) bereits vollständig enthält,
beschränkte sich die Multi-GPU-Aktivierung auf die \emph{Verkabelung}: Der
Trainingsstarter wurde von \texttt{python} auf
\texttt{torchrun -{}-nproc\_per\_node=}\allowbreak\texttt{\$NUM\_GPUS} umgestellt (bei
\texttt{NUM\_GPUS=1} bleibt es rückwärtskompatibel beim einfachen
\texttt{python}), und die SLURM-Ressourcen wurden angepasst.

\begin{warningbox}[Host-RAM-Engpass bei verteilten Dataloadern]
  Eine nicht offensichtliche Hürde war der Arbeitsspeicher des Hosts: Bei
  acht Dataloader-Prozessen pro Rang und vier GPUs entstehen 32~Prozesse, die
  256~GB RAM überliefen und vom System beendet wurden
  (\texttt{DataLoader worker killed by signal: Killed}). Die Lösung bestand darin,
  \texttt{DATALOADER\_WORKERS=4} zu setzen und den angeforderten Speicher auf
  384~GB zu erhöhen. Details: \doc{docs/training/multi-gpu.md}.
\end{warningbox}

Erwartet wird durch NVLink eine Beschleunigung um etwa das 3,5-fache (nicht das
volle Vierfache, da die NCCL-Kommunikation 5--20\,\% kostet). Der reale
Vier-GPU-Lauf wurde im Rahmen dieser Arbeit nicht zeitlich vermessen; der
ausgewertete vollständige Lauf ist der Einzel-GPU-Lauf (\cref{sec:trainingsverlauf}).

\section{Durchführung auf dem KISSKI-Cluster}
\label{sec:kisski}

Das Training lief auf dem \ac{hpc}-Cluster \textbf{KISSKI} der GWDG
Göttingen~\cite{kisski}, der zwei GPU-Partitionen bereitstellt: \texttt{kisski}
(NVIDIA~A100, \qty{80}{\giga\byte}) und \texttt{kisski-h100} (H100,
\qty{94}{\giga\byte}), jeweils mit maximal 48~h Walltime. Jobs werden per
\ac{slurm}-Batch-Skript eingereicht (Auszug in \cref{anh:slurm}); die Laufzeit
nutzt Apptainer statt Docker. Da die Rechenknoten keinen Internetzugang besitzen,
ergaben sich zwei praktische Besonderheiten:

\begin{itemize}
  \item \textbf{W\&B im Offline-Modus.} \ac{wandb} läuft zwangsweise offline; die
    Logs werden nachträglich vom Login-Knoten synchronisiert
    (\doc{docs/training/wandb-offline-sync.md}).
  \item \textbf{Variablenübergabe.} Geheimnisse wie \texttt{HF\_TOKEN} müssen als
    Inline-Präfix vor \texttt{sbatch} bzw.\ über eine Schlüsseldatei übergeben
    werden, da KISSKI die Umgebungsweitergabe einschränkt
    (\doc{docs/training/kisski-hpc.md}).
\end{itemize}

\section{Evaluationsmethodik}
\label{sec:eval-methodik}

Wie in \cref{sec:domain-gap-grundlagen} eingeführt, wurde zweigleisig evaluiert.

\paragraph{Open-Loop-Action-MSE.}
Dies ist NVIDIAs offizielle Checkpoint-Metrik und der Goldstandard für die reine
\emph{Modellqualität}: Der \ac{mse} zwischen vorhergesagter und demonstrierter
Aktion wird auf gehaltenen echten Trajektorien gemessen, ohne Regelkreis und
damit ohne Domain-Gap. Ein absoluter Schwellwert für \enquote{gut} existiert
nicht; die Metrik dient dem relativen Vergleich von Checkpoints.

\paragraph{Checkpoint-Auswahl durch nachgelagerten Sweep.}
Aus der fehlenden Validierung während des Trainings folgt ein eigener
Arbeitsschritt. Die Einschränkung ist dabei nicht bloß Konfiguration, sondern
\emph{strukturell}: Der verwendete Fork erzwingt \texttt{eval\_strategy="no"}
per Assertion, die vorhandenen Konfigurationsfelder für eine Open-Loop-Eval
während des Trainings werden nie ausgelesen. Die Auswahl des besten Checkpoints
muss daher \emph{nach} dem Lauf erfolgen. Dafür misst
\doc{Training/scripts/checkpoint\_sweep.py} den Open-Loop-\ac{mse} \emph{jedes}
gespeicherten Checkpoints auf den zurückgehaltenen Testepisoden und wählt das
Minimum. Dieser Schritt ersetzt die zuvor mangels Alternative geübte Praxis,
schlicht den letzten Checkpoint zu nehmen -- eine Praxis, die sich im Nachhinein
als falsch erwies (\cref{sec:lauf3}).

\paragraph{Closed-Loop-Sim-Eval.}
Für die geschlossene Schleife wurde eine Isaac-Lab-Umgebung mit Roboter, Tisch,
Würfeln und fünf Kameras aufgebaut. Architektonisch sind zwei Prozesse über
\ac{zmq} gekoppelt (\cref{fig:zmq}): ein \textbf{GR00T-Policy-Server}, der den
Checkpoint lädt und Aktionsanfragen beantwortet, und ein \textbf{Isaac-Lab-Sim-Client},
der Bilder rendert und Aktionen ausführt. Die Trennung ist nötig, weil beide
Komponenten inkompatible Python-Umgebungen verwenden (Server: Python~3.10;
Isaac~Sim: gebündeltes Python~3.11). Die vollständige Bedienung beschreibt
\doc{docs/simulation/vastai-anleitung.md}.

\begin{figure}[H]
  \centering
  \begin{tikzpicture}[
    node distance = 2.2cm,
    proc/.style  = {draw, rounded corners, fill=NavyBlue!8,
                    minimum width=3.6cm, minimum height=1.4cm,
                    font=\small\sffamily, align=center},
    arrow/.style = {-{Stealth[scale=1.0]}, thick, NavyBlue!70},
  ]
    \node[proc] (server) {GR00T-Policy-Server\\\footnotesize Python 3.10, \textasciitilde10\,GB VRAM\\\footnotesize lädt \texttt{checkpoint-*}};
    \node[proc] (client) [right=3.2cm of server] {Isaac-Lab-Sim-Client\\\footnotesize Python 3.11, \textasciitilde8\,GB VRAM\\\footnotesize rendert 4+1 Kameras};
    \draw[arrow] (client.north) to[bend right=22]
      node[above, font=\scriptsize\sffamily] {Beobachtung (REQ)} (server.north);
    \draw[arrow] (server.south) to[bend right=22]
      node[below, font=\scriptsize\sffamily] {Aktion (REP)} (client.south);
  \end{tikzpicture}
  \caption{Client-Server-Architektur der Closed-Loop-Evaluation. Sim-Client und
    Policy-Server kommunizieren über \ac{zmq} (Port 5555) auf demselben Host.}
  \label{fig:zmq}
\end{figure}

\begin{warningbox}[Hardware-Anforderung: RT-Cores]
  Die Closed-Loop-Eval stellt eine ungewöhnliche Hardware-Anforderung: Isaac~Sim
  rendert die Kamerabilder per Raytracing und benötigt dafür \textbf{RT-Cores}
  (Ampere-Generation oder neuer, z.\,B.\ L40, RTX~4090, A6000).
  \emph{Ausgerechnet die starken Trainings-GPUs A100 und H100 besitzen keine
  RT-Cores} und sind daher für die Sim-Eval ungeeignet -- die
  KISSKI-Trainingspartitionen können die Closed-Loop-Eval nicht ausführen. Die
  Sim-Eval lief deshalb zunächst auf vast.ai (NVIDIA~L40); ab Juli~2026 stand
  ein institutseigener Server mit zwei RT-Core-fähigen RTX~PRO~6000
  (Blackwell) zur Verfügung, auf dem auch die Referenz-Eval und das
  RL-Fine-Tuning liefen. Hintergrund:
  \doc{docs/simulation/umsetzungsnotizen.md}.
\end{warningbox}

\paragraph{Inferenz-Durchsatz.}
Da jede Episode viele hundert Modellaufrufe erfordert, wurde die Inferenzseite
gesondert vermessen und optimiert: Der Flow-Matching-Aktionskopf lässt sich in
ein GPU-spezifisches Laufzeitformat übersetzen, und das Rendern der Kameras
lässt sich auf den Takt der Politikaufrufe synchronisieren, statt jedes
Simulationsbild zu rendern. Methodisch bemerkenswert ist dabei die Trennung der
Messgrößen: Die reine \emph{Politik}-Latenz wird ohne Simulation und ohne
Netzwerktransport bestimmt und ist damit die einzige Durchsatzgröße dieser
Arbeit, die auch auf realer Hardware gilt -- das Rendern, das in der Simulation
den Löwenanteil der Laufzeit ausmacht, entfällt dort ersatzlos.
Details: \doc{docs/simulation/inferenz-optimierung.md}.

\paragraph{Referenz-Eval zur Pipeline-Validierung.}
Ergänzend zu den beiden internen Evaluationsgleisen wurde eine \emph{externe}
Kontrolle vorgesehen: Das un-feinabgestimmte Basismodell wird auf einer
Referenzaufgabe mit von NVIDIA publizierten Erfolgsquoten (RoboCasa GR-1
Tabletop) ausgeführt. Reproduziert die eigene Inferenz-Kette diese Sollwerte,
sind Checkpoint-Laden, Embodiment-Konfiguration und Observation-/Action-Mapping
als Fehlerquellen für ein Null-Resultat ausgeschlossen. Durchführung und
Ergebnis beschreibt \cref{sec:robocasa}.

\paragraph{Baseline.}
Als unterer Vergleichswert war eine Baseline mit dem \emph{un}-feinabgestimmten
Modell auf dem Standard-G1-Embodiment vorgesehen. Da dieses Embodiment auf
Ganzkörper-Loko-Manipulation mit nur einer Kamera ausgelegt ist und stark
außerhalb der Verteilung der Tabletop-Szene liegt, ist hier eine Erfolgsrate
nahe~$0\,\%$ zu erwarten (\doc{docs/simulation/baseline-eval.md}).
